\documentclass[../../main.tex]{subfiles}
\graphicspath{{../../image/}{../../figure/}}
\begin{document}

\section{广义特征线}\label{section:ch10-广义特征线}

与古典的、Lipschitz 连续的解相联系的特征线，是在\cref{chapter:ch07-一维双曲平衡律组}中通过\cref{definition:ch07-i特征线}引入的。它们是古典理论用于研究解的分析性质与几何性质的主要工具之一。因此，把这一概念推广到弱解框架是十分自然的。


此处我们将与弱解 $U$ 相联系的 $i$-特征族的特征线，完全按古典情形那样定义为常微分方程\cref{chapter:ch07-一维双曲平衡律组}在 Filippov 意义下的积分曲线：


\begin{definition}[广义 $i$-特征线]\label{definition:ch10-广义i特征线}
对于方程组\eqref{align:ch10-10-1-1}，与（一般为弱的）解 $U$ 相联系、定义在时间区间 $[\sigma,\tau]\subset[0,\infty)$ 上的一条\emph{广义 $i$-特征线}，是指一个 Lipschitz 函数 $\xi:[\sigma,\tau]\to(-\infty,\infty)$，它满足微分包含
\begin{align}
\dot{\xi}(t)\in\Lambda_{i}(\xi(t),t),\qquad\text{a.e. on }[\sigma,\tau],
\label{align:ch10-10-2-1}
\end{align}
其中
\begin{align}
\Lambda_{i}(\bar{x},\bar{t})\coloneqq\bigcap_{\varepsilon>0}\left[\operatorname*{ess\,inf}_{[\bar{x}-\varepsilon,\bar{x}+\varepsilon]}\lambda_{i}(U(x,\bar{t})),\ \operatorname*{ess\,sup}_{[\bar{x}-\varepsilon,\bar{x}+\varepsilon]}\lambda_{i}(U(x,\bar{t}))\right].
\label{align:ch10-10-2-2}
\end{align}
\end{definition}

由形如\eqref{align:ch10-10-2-1}的偶然方程的一般理论，立即得到下述结论：

\begin{theorem}\label{theorem:ch10-过点两条广义特征线}
过任意固定点 $(\bar{x},\bar{t})\in(-\infty,\infty)\times[0,\infty)$，存在两条（不必互异的）与 $U$ 相联系、定义在 $[0,\infty)$ 上的广义 $i$-特征线，即极小者 $\xi_{-}(\cdot)$ 与极大者 $\xi_{+}(\cdot)$，满足 $\xi_{-}(t)\leqslant\xi_{+}(t)$，$t\in[0,\infty)$。夹在 $\xi_{-}(\cdot)$ 与 $\xi_{+}(\cdot)$ 的图像之间的漏斗形区域，由所有可用一条与 $U$ 相联系的广义 $i$-特征线与 $(\bar{x},\bar{t})$ 相连通的点 $(x,t)$ 所组成。
\end{theorem}

微分包含解的其他标准性质对广义特征线理论也有有用的推论：若 $\{\xi_{m}(\cdot)\}$ 是一列与 $U$ 相联系、定义在 $[\sigma,\tau]$ 上的广义 $i$-特征线，在 $[\sigma,\tau]$ 上一致收敛于某个 Lipschitz 函数 $\xi(\cdot)$，则 $\xi(\cdot)$ 必为一条与 $U$ 相联系的广义 $i$-特征线。特别地，若 $\xi_{m}(\cdot)$ 是过点 $(x_{m},\bar{t})$ 的极小（或极大）广义 $i$-特征线，且 $x_{m}\downarrow\bar{x}$（相应地 $x_{m}\uparrow\bar{x}$），则当 $m\to\infty$ 时 $\{\xi_{m}(\cdot)\}$ 收敛于过点 $(\bar{x},\bar{t})$ 的极小（相应地极大）广义 $i$-特征线 $\xi_{-}(\cdot)$（相应地 $\xi_{+}(\cdot)$）。

除古典 $i$-特征线外，满足 Lax $E$ 条件的 $i$-激波也是广义 $i$-特征线的显然实例。事实上，可以证明它们是仅有的两类可能情形。的确，尽管\cref{definition:ch10-广义i特征线}表面上允许 $\dot{\xi}$ 在区间 $\Lambda_{i}$ 中取任意值，但 $U$ 是\eqref{align:ch10-10-1-1}的解这一事实，约束了与 $U$ 相联系的广义 $i$-特征线只能以古典 $i$-特征速度或以 $i$-激波速度传播：

\begin{theorem}\label{theorem:ch10-广义特征线速度约束}
设 $\xi(\cdot)$ 是一条与 $U$ 相联系、定义在 $[\sigma,\tau]$ 上的广义 $i$-特征线。则对几乎所有 $t\in[\sigma,\tau]$，有如下结论：当 $(\xi(t),t)\in\mathscr{C}$ 时，$\dot{\xi}(t)=\lambda_{i}(U_{0})$，其中 $U_{0}=U(\xi(t)\pm,t)$；当 $(\xi(t),t)\in\mathscr{J}$ 时，$\dot{\xi}(t)=s$，其中 $s$ 是左侧状态 $U_{-}$ 与右侧状态 $U_{+}$ 之间 $i$-激波的速度，$U_{\pm}=U(\xi(t)\pm,t)$。特别地，$s$ 既满足 Rankine--Hugoniot 条件\cref{chapter:ch08-容许激波}，也满足 Lax $E$ 条件\eqref{align:ch10-10-1-2}。

\end{theorem}

\begin{proof}
回顾\cref{section:ch10-bv解}中记录的 $BV$ 函数性质便知：对几乎所有使 $(\xi(t),t)\in\mathscr{C}$ 的 $t\in[\sigma,\tau]$，区间 $\Lambda_{i}(\xi(t),t)$ 退化为单点 $\lambda_{i}(U(\xi(t)\pm,t))$，从而由\eqref{align:ch10-10-2-1}得 $\dot{\xi}(t)=\lambda_{i}(U(\xi(t)\pm,t))$。

将测度等式\eqref{align:ch10-10-1-1}应用于 $\xi$ 图像的任意子弧，并利用\cref{theorem:ch01-Gauss-Green}（特别是其中的\eqref{align:ch01-GG-jump}），得到
\begin{align}
F(U(\xi(t)+,t))-F(U(\xi(t)-,t))=\dot{\xi}(t)\,[U(\xi(t)+,t)-U(\xi(t)-,t)],
\label{align:ch10-10-2-3}
\end{align}
在 $[\sigma,\tau]$ 上几乎处处成立。因此，对几乎所有使 $(\xi(t),t)\in\mathscr{J}$ 的 $t\in[\sigma,\tau]$，有 $\dot{\xi}(t)=s$，其中 $s$ 是连接两状态 $U_{-}=U(\xi(t)-,t)$ 与 $U_{+}=U(\xi(t)+,t)$ 的激波速度。由解的结构，存在 $j\in\{1,\cdots,n\}$ 使得
\begin{align*}
\lambda_{j-1}(U_{\pm})<\lambda_{j}(U_{+})\leqslant s\leqslant\lambda_{j}(U_{-})<\lambda_{j+1}(U_{\pm}).
\end{align*}
另一方面，由\eqref{align:ch10-10-2-1}可知 $s$ 落在以 $\lambda_{i}(U_{-})$ 与 $\lambda_{i}(U_{+})$ 为端点的区间内。因此 $j=i$，且\eqref{align:ch10-10-1-2}成立。证明完毕。

\end{proof}

上述定理引出如下术语：

\begin{definition}[无激波广义 $i$-特征线]\label{definition:ch10-无激波广义i特征线}
一条与 $U$ 相联系、定义在 $[\sigma,\tau]$ 上的广义 $i$-特征线 $\xi(\cdot)$，若对几乎所有 $t\in[\sigma,\tau]$ 都有 $U(\xi(t)-,t)=U(\xi(t)+,t)$，则称为\emph{无激波的}。
\end{definition}

由\cref{theorem:ch10-广义特征线速度约束}的证明可知，\eqref{align:ch10-10-2-1}等价于
\begin{align}
\dot{\xi}(t)\in\left[\lambda_{i}(U(\xi(t)+,t)),\ \lambda_{i}(U(\xi(t)-,t))\right],\qquad\text{a.e. on }[\sigma,\tau].
\label{align:ch10-10-2-4}
\end{align}

在下文，一类以最大或最小容许速度传播的特殊广义特征线将扮演重要角色：

\begin{definition}[左/右 $i$-接触]\label{definition:ch10-左右i接触}
一条与 $U$ 相联系、定义在 $[\sigma,\tau]$ 上的广义 $i$-特征线 $\xi(\cdot)$，若满足
\begin{align}
\dot{\xi}(t)=\lambda_{i}(U(\xi(t)-,t)),\qquad\text{a.e. on }[\sigma,\tau],
\label{align:ch10-10-2-5}
\end{align}
则称为\emph{左 $i$-接触}；若满足
\begin{align}
\dot{\xi}(t)=\lambda_{i}(U(\xi(t)+,t)),\qquad\text{a.e. on }[\sigma,\tau],
\label{align:ch10-10-2-6}
\end{align}
则称为\emph{右 $i$-接触}。
\end{definition}

显然，无激波 $i$-特征线同时是左、右 $i$-接触。注意，由于它们是广义 $i$-特征线，左（相应地右）$i$-接触也应满足\cref{theorem:ch10-广义特征线速度约束}的结论，即对几乎所有使 $(\xi(t),t)\in\mathscr{J}$ 的 $t\in[\sigma,\tau]$ 有 $\dot{\xi}(t)=s$。当然，在不容许左（相应地右）接触间断的方程组中这是不可能的。在这样的方程组中，左（相应地右）接触必然是无激波的。特别地，回顾\cref{chapter:ch08-容许激波}可知：当方程组\eqref{align:ch10-10-1-1}的 $i$-特征族真正非线性且 $U$ 的振幅充分小时，任何左或右 $i$-接触必为无激波的。


\end{document}